\chapter{Conclusions and Recommendations}
\label{ch:conclusions_and_recommendations}

This chapter synthesizes the primary empirical findings, theoretical insights, and operational conclusions derived from the two-stage multi-commodity Mixed-Integer Linear Programming (MILP) model developed for post-disaster resource allocation in Davao City. It outlines the specific scholarly contributions of this research, articulates actionable policy recommendations for municipal disaster management agencies, candidly discusses the methodological limitations of the study, and identifies promising avenues for future investigation in humanitarian operations research.

% ==============================================================================
% 5.1 OVERVIEW AND SUMMARY OF FINDINGS
% ==============================================================================
\section{Overview and Summary of Findings}
\label{sec:summary_of_findings}

Natural disasters inflict severe logistical strain on municipal emergency response systems. In large, topographically complex jurisdictions like Davao City---spanning 182 administrative barangays encompassing dense coastal urban centers, sprawling river valleys, and isolated mountainous highlands---the fundamental challenge of humanitarian logistics is allocating acutely scarce relief commodities in a manner that is simultaneously computationally tractible, physiologically grounded, and distributively equitable. 

Conventional emergency distribution systems in local government units (LGUs) frequently succumb to two well-documented systemic failure modes:
\begin{enumerate}
    \item \textbf{Decentralized Order Amplification (The Bullwhip Effect):} Uncoordinated relief requests generated by local community leaders under crisis panic lead to extreme order distortion, resulting in over-allocation and stockpiling at accessible urban hubs while leaving remote and critically affected populations without life-sustaining aid.
    \item \textbf{The Utilitarian Small-Demand Paradox:} Standard unweighted satisfaction metrics ($\sum \phi_{jr}$) create strong algorithmic incentives to divert relief away from high-density, high-hazard urban centers to small rural barangays, where minimal supply volumes can drive fulfillment ratios to 100\%, thereby maximizing superficial percentage statistics while exacerbating aggregate human mortality.
\end{enumerate}

To resolve these operational and ethical dilemmas, this study developed an exact two-stage multi-commodity Mixed-Integer Linear Programming formulation that couples discrete facility activation decisions across municipal logistics hubs ($H_{\max} = 4$ of 6 candidate Disaster Risk Reduction and Management Operations Centers) with continuous resource dispatch flows across two survival-critical commodities: potable water and emergency food calories. By replacing arbitrary volumetric measures with non-linear exponential deprivation cost functions ($e^{\lambda_r(1-\phi_{jr})} - 1$) linearized through a certified 20-segment piecewise outer tangent epigraph formulation, the model explicitly accounts for the non-linear, time-dependent escalation of human suffering resulting from commodity deficits. Concurrently, the framework strictly enforces universal anti-starvation survival floors ($\varepsilon_{\min} = 0.10$) and anti-hoarding demand ceilings ($\phi_{jr} \le 1.0$) across all 182 barangays.

The empirical evaluation conducted in Chapter~4 established seven major findings:
\begin{itemize}
    \item \textbf{Sub-Second Computational Tractability:} On standard local workstation hardware (Intel Core i5, 8~GB RAM), the exact model solves to certified global optimality ($0.000\%$ relative MIP gap) in an average of \textbf{0.696 seconds} ($T_{\text{solve}}$) with a memory footprint of \textbf{118.4 MB}, requiring \textbf{zero branch-and-bound tree exploration} due to root-node linear relaxation convergence.
    \item \textbf{High-Fidelity Linearization:} The 20-segment outer tangent discretization bounds maximum absolute epigraph under-approximation error to $\Delta_{\max} \le 0.005890$ ($< 1.99\%$ for potable water, $< 0.50\%$ for food), preserving over $98.0\%$ to $99.5\%$ mathematical fidelity to the true exponential deprivation curves without incurring the computational intractability of non-linear solvers.
    \item \textbf{Superiority Over Linear and Proportional Heuristics:} The proposed model achieves a global minimum deprivation cost of $Z^* = \mathbf{2.8283}$, reducing human suffering by \textbf{11.32\%} relative to uniform proportional and egalitarian allocation ($Z = 3.1893$) and by \textbf{1.66\%} relative to a linear hazard-weighted proportional heuristic ($Z = 2.8761$).
    \item \textbf{Transparent Characterization of Equity Trade-Offs:} An objective multi-criteria evaluation demonstrates that the proposed model does not dominate on neutral equity metrics: the Proportional and Egalitarian baselines achieve higher floor fulfillment ($\phi_{\min} = 0.3333$ vs. $0.1000$), zero floor-bound population share ($0.0\%$ vs. $56.0\%$), and higher Nash social welfare ($W_{\text{Nash}} = -199.95$ vs. $-359.13$). The proposed model intentionally trades off floor parity and Nash welfare to aggressively concentrate scarce relief into severe-hazard zones ($\bar{\phi}_{\text{High}} = 0.4821$, hazard parity ratio $4.821$).
    \item \textbf{Automated Policy Diagnostic Boundaries:} Scenario stress-testing across seven calamity profiles identified a sharp systemic insolvency threshold at effective inventory levels below 10\% (SC-04), where the model reports infeasibility, serving as an automated decision-support trigger for immediate Tier-3 National Escalation. Furthermore, identical 60\% deficits produce a 60.8\% higher deprivation penalty for potable water than for food ($Z^* = 4.85$ vs. $3.02$), validating physiological urgency prioritization.
    \item \textbf{Mitigation of Accessibility Bias and Order Amplification:} The centralized MILP suppresses road accessibility correlation to $r = +0.2891$ (compared to $r = +1.0000$ in accessibility-driven distribution) and prevents uncoordinated over-allocation by mathematical construction through the demand ceiling ($0$ wasted units vs. $2.6$ million wasted units in decentralized Monte Carlo simulations).
    \item \textbf{Rigorous Non-Parametric Validation:} Omnibus Friedman ($\chi_F^2 = 384.21, p < 0.0001$), post-hoc Nemenyi ($\text{CD} = 0.4518$), and paired Wilcoxon signed-rank tests across $n = 182$ paired barangay allocations within scenarios ($Z \ge 8.42, p < 10^{-6}$) confirm that all observed performance gains are statistically significant.
\end{itemize}


% ==============================================================================
% 5.2 CONCLUSIONS
% ==============================================================================
\section{Conclusions}
\label{sec:conclusions}

Based on the empirical findings, mathematical verifications, and trade-off analyses, the following primary conclusions are drawn:

\subsection{Algorithmic Tractability and Real-Time Operational Viability}
The proposed Two-Stage MILP resolves the computational barrier long associated with combinatorial disaster logistics models. With total execution time (model generation plus solve time) remaining strictly below 0.80 seconds across the entire 182-barangay municipal network of Davao City, the framework is proven to be computationally viable for real-time emergency operations center (EOC) dispatch. The fact that the model converges at the root node without requiring branch-and-bound tree enumeration confirms that the underlying polyhedral structure of the two-stage hub-and-spoke allocation problem possesses favorable integrality properties under proportional network capacities. Consequently, local governments do not require supercomputing clusters, cloud infrastructure, or commercial solver licenses; the model operates reliably using open-source solvers (such as HiGHS) on standard field laptops.

\subsection{Mathematical Fidelity of Piecewise Tangent Approximation}
Piecewise linear outer tangent epigraph formulation provides an exact, robust, and computationally stable methodology for handling non-linear human suffering functions in humanitarian logistics. By selecting $K = 20$ linear segments ($21$ breakpoints at uniform $\Delta \phi = 0.05$ intervals), the formulation bounds relative approximation error to under 2.0\% for water and under 0.5\% for food. This level of mathematical fidelity eliminates the risk of numerical instability, local optima trapping, and non-convex convergence failures inherent in general non-linear programming (NLP) solvers, while maintaining complete compatibility with standard linear branch-and-cut engines.

\subsection{Humanitarian Deprivation Minimization vs. Linear Heuristics}
Exact mathematical optimization provides tangible, quantifiable life-saving value beyond what intuitive operational heuristics can achieve. While simple linear hazard-weighting heuristics improve upon unweighted proportional distribution by prioritizing vulnerable communities, they cannot account for the exponential rate at which human physiological harm escalates as fulfillment declines. The proposed MILP actively traces the steepest gradient of the convex deprivation surface, systematically reallocating marginal supplies where the derivative of suffering is greatest. This non-linear optimization unlocks an additional 1.66\% reduction in total human suffering beyond hazard weighting alone, and an 11.32\% reduction compared to status-quo proportional distribution, directly demonstrating the necessity of mathematical programming over static spreadsheet heuristics.

\subsection{Distributive Justice, Equity Trade-Offs, and Reconciliation with ``Fairness to All''}
A central philosophical contribution of this research is clarifying the operational meaning of distributive equity in resource-constrained disaster relief. Under acute supply scarcity (e.g., municipal inventories covering only 33.3\% of aggregate daily need), classical egalitarianism ($\phi_j = 0.3333$ for all $j$) and Nash welfare maximization represent equality of process rather than equity of outcome. Treating an upland community that is physically unaffected by flooding identically to a submerged riverbank community where infants and elderly individuals face rapid dehydration is ethically indefensible in life-or-death crisis management.

The proposed framework successfully reconciles distributive justice through a two-tier humanitarian compact:
\begin{enumerate}
    \item \textbf{Guaranteed Universal Baseline:} Every single community, regardless of geographic remoteness or population density, is guaranteed a non-negotiable survival allocation ($\phi_{jr} \ge \varepsilon_{\min} = 0.10$), ensuring that peripheral upland barangays (such as Marilog and Paquibato) are never abandoned in pursuit of aggregate efficiency.
    \item \textbf{Need-Prioritized Life-Saving Surplus:} All discretionary relief stock above the 10\% survival floor is channeled dynamically to communities facing catastrophic hazard exposure, raising their fulfillment ratios to $48\%$--$57\%$.
\end{enumerate}
Thus, the model's intentional sacrifice of floor parity (leaving 56\% of the population at the 10\% floor) and lower Nash social welfare is not an operational failure; it is the mathematical hallmark of a rational, prioritized harm-reduction doctrine designed to prevent acute mortality.

\subsection{Systemic Insolvency Boundaries and Automated Crisis Diagnostics}
Mathematical model infeasibility, when properly structured, serves as a vital diagnostic mechanism rather than a computational failure. When total available active warehouse inventory falls below the universal 10\% survival threshold (as demonstrated in SC-04 at $\alpha = 0.25$, yielding 8.33\% effective throughput), the solver's declaration of infeasibility provides an automated, indisputable mathematical certification that municipal resources are insolvent. This boundary diagnostic prevents local officials from engaging in homeopathic aid dilution (distributing therapeutically useless fractions of aid to project false coverage) and serves as an automated trigger for immediate Tier-3 National Escalation.

\subsection{Spatial Equity and Mitigation of Order Distortion}
Centralized optimization effectively neutralizes the spatial biases and informational dysfunctions that plague decentralized disaster response. In uncoordinated field operations, accessible urban barangays close to central depots aggressively inflate relief requests, generating severe order amplification (bullwhip distortion) that wastes millions of relief units while starving peripheral and isolated communities. The proposed MILP prevents over-allocation by mathematical construction ($0$ units wasted) through the unified demand ceiling ($\phi_{jr} \le 1.0$), while simultaneously reducing spatial accessibility bias from a perfect correlation ($r = +1.0000$ by benchmark construction) down to a weak $r = +0.2891$.



% ==============================================================================
% 5.3 THEORETICAL AND METHODOLOGICAL CONTRIBUTIONS
% ==============================================================================
\section{Theoretical and Methodological Contributions}
\label{sec:contributions}

This research advances the literature on humanitarian operations research, mathematical programming, and disaster management through three core theoretical and methodological contributions:

\paragraph{1. Multi-Commodity Deprivation Formulation with Post-Hoc Liebig Integration:}
Prior literature on humanitarian deprivation costs (e.g., Holgu{\'i}n-Veras et al., 2013; P{\'e}rez-Rodr{\'i}guez and Holgu{\'i}n-Veras, 2016) has predominantly evaluated single-commodity allocations or unweighted multi-item aggregates. This study extends the formal Two-Stage MILP structure to simultaneously track distinct survival commodities (water and food) with physiologically differentiated deprivation escalation parameters ($\lambda_{\text{water}} = 2.0$ vs. $\lambda_{\text{food}} = 0.5$). Furthermore, by enforcing anti-starvation constraints directly per resource within the model ($\phi_{jr} \ge \varepsilon_{\min}$) and evaluating overall community fulfillment post-hoc via Liebig's Law of the Minimum ($\phi_j = \min_r \phi_{jr}$), the formulation eliminates the underdetermination and solver indifference bugs that afflict naive multi-commodity LP models where aggregate fulfillment is uncoupled from the objective function.

\paragraph{2. Zero-Baseline Normalized Shift Formulation:}
In standard convex deprivation formulations, the cost function $e^{\lambda(1-\phi)}$ evaluates to $1.0$ at full satisfaction ($\phi = 1.0$), introducing an arbitrary constant baseline penalty into the objective function that distorts cross-model comparisons. This study rigorously establishes and linearizes the zero-baseline shifted deprivation formulation:
\begin{equation}
    g_r(\phi_{jr}) = e^{\lambda_r(1-\phi_{jr})} - 1
\end{equation}
ensuring that complete commodity satisfaction ($\phi_{jr} = 1.0$) yields identically zero deprivation cost ($g_r(1) = 0$). This ensures mathematical consistency across different scale factors and enables valid comparative evaluations against utilitarian, egalitarian, and proportional baselines.

\paragraph{3. Decoupled Multi-Criteria Equity Evaluation Suite:}
A persistent methodological flaw in humanitarian logistics research is evaluating models solely against their own primary objective function (circular validation). This study designs and executes an independent, multi-criteria distributive equity benchmark comprising worst-off fulfillment ($\phi_{\min}$), floor-bound population share ($f_{\text{floor}}$), fulfillment-ratio variance ($\sigma_\phi$), Nash social welfare ($W_{\text{Nash}}$), fulfillment disparities ($\Delta_\varepsilon$), and hazard parity ratios ($\bar{\phi}_{\text{High}} / \bar{\phi}_{\text{Min}}$). This decoupled evaluation framework allows researchers and policymakers to quantify the exact trade-offs between efficiency, egalitarian equality, and deprivation minimization.


% ==============================================================================
% 5.4 OPERATIONAL AND POLICY RECOMMENDATIONS
% ==============================================================================
\section{Policy and Operational Recommendations for Davao City DRRMOC}
\label{sec:policy_recommendations}

The empirical findings and operational analyses of this study translate directly into concrete, actionable policy recommendations for the City Disaster Risk Reduction and Management Office (CDRRMO), the City Social Welfare and Development Office (CSWDO), and the Department of Social Welfare and Development (DSWD Field Office XI):

\begin{enumerate}
    \item \textbf{Institutionalize the Non-Negotiable 10\% Survival Baseline ($\varepsilon_{\min} = 0.10$):}
    Municipal disaster contingency guidelines should formally mandate that no barangay within Davao City receive less than 10\% of its daily Sphere-standard survival allocation ($1.5\text{ L}$ water and $210\text{ kcal}$ food per person/day) during the first 24 to 48 hours of an acute crisis. Enforcing this universal survival floor in municipal relief protocols guarantees that remote peripheral barangays (such as Paquibato and Marilog) cannot be bypassed or deprioritized due to transportation distance or communication blackouts.

    \item \textbf{Adopt Algorithmic Diagnostic Boundaries for Tier-3 National Escalation:}
    Municipal disaster protocols should establish an automated inventory solvency threshold. When available prepositioned warehouse inventory drops below 10\% of aggregate city-wide daily demand (as identified in SC-04), the CDRRMO should immediately bypass localized triage and initiate an automated Tier-3 National Escalation request to the National Disaster Risk Reduction and Management Council (NDRRMC) and DSWD Central Office for regional airlifts and inter-regional supply augmentation.

    \item \textbf{Rebalance Prepositioned Warehousing Toward Potable Water Reserves:}
    Given that acute water deficits generate a 60.8\% higher human deprivation penalty than identical caloric deficits ($Z_{\text{water}} = 4.85$ vs. $Z_{\text{food}} = 3.02$), municipal disaster budgets should rebalance contingency stockpiles. Rather than predominantly investing in non-perishable dry food rations (family food packs), the City of Davao should prioritize mobile water treatment units, chemical purification tablets, modular bladder tanks, and solar-powered filtration stations prepositioned across its six strategic logistics hubs.

    \item \textbf{Transition from Decentralized Request Queues to Centralized Single-Source Optimization:}
    The city should phase out its historical reliance on uncoordinated barangay relief requests processed on a first-come, first-served basis. The CDRRMO should deploy the proposed Two-Stage MILP as the central algorithmic engine of its Disaster Operations Center, establishing a single mathematical source of truth that prevents panic-driven request inflation, eliminates multi-million-unit relief waste, and ensures equitable geographic distribution.

    \item \textbf{Integrate High-Resolution Geohazard GIS Layers into Disaster Allocation Databases:}
    The Davao City Government should formally link its Philippine Statistics Authority demographic registries (POPCEN 2024) and PSGC codes with digitized Mines and Geosciences Bureau (MGB) flood and landslide susceptibility layers. Automatically embedding real-time hazard severity indices ($H_j$) into barangay priority weights ($P_j$) ensures that the optimization model can be populated and executed within minutes of disaster onset, replacing ad-hoc manual estimations with empirical hazard data.
\end{enumerate}


% ==============================================================================
% 5.5 LIMITATIONS OF THE STUDY
% ==============================================================================
\section{Limitations of the Study}
\label{sec:limitations}

To maintain rigorous academic honesty and provide appropriate context for interpreting the results, the methodological and empirical limitations of this study must be explicitly recognized:

\begin{enumerate}
    \item \textbf{Deterministic Decision Structure Without Stochastic Recourse:}
    The formulation models hub activation (Stage~1) and flow allocation (Stage~2) as a deterministic two-stage problem optimized simultaneously under complete information. It does not represent classical two-stage stochastic programming with recourse, wherein first-stage decisions are made under probabilistic uncertainty regarding weather conditions, population displacement, or secondary hazard impacts.

    \item \textbf{Exclusion of Detailed Vehicle Routing and Network Damage Dynamics:}
    The current model focuses on the macro-allocation problem (determining how much aggregate supply should flow from specific hubs to specific barangays) and assumes direct arc connectivity. It does not solve the underlying Vehicle Routing Problem (VRP) or model micro-level road network disruptions, bridge collapses, flood-inundated transit corridors, or vehicle fleet dispatch schedules.

    \item \textbf{Provisionality of Physiological Deprivation Severity Parameters:}
    While the relative ordering of deprivation rates ($\lambda_{\text{water}} = 2.0 > \lambda_{\text{food}} = 0.5$) is grounded in clinical medicine and the physiological ``Rule of Threes,'' the exact numerical magnitudes represent provisional calibration values rather than empirically measured willingness-to-pay or mortality-rate metrics derived from local post-disaster epidemiological surveys.

    \item \textbf{Single-Period Static Dispatch Horizon:}
    The evaluation models an acute 24-hour post-disaster dispatch window. It does not model multi-period rolling-horizon dynamics, inventory replenishment from outside regions over multi-week horizons, or the cumulative compounding of deprivation penalties across sequential operational cycles.

    \item \textbf{Administrative Boundary Demarcation Limitations:}
    Demand aggregation is conducted at the administrative barangay level ($n = 182$) based on census data. This assumes uniform internal vulnerability within each barangay, which may obscure intra-barangay socio-economic disparities and localized micro-topographical flooding variations within large territorial barangays.
\end{enumerate}


% ==============================================================================
% 5.6 DIRECTIONS FOR FUTURE RESEARCH
% ==============================================================================
\section{Directions for Future Research}
\label{sec:future_research}

Building upon the theoretical foundation and empirical findings established in this thesis, several promising directions for future research are identified:

\begin{enumerate}
    \item \textbf{Two-Stage Stochastic Programming with Recourse:}
    A natural and vital theoretical extension is formulating the problem as a true Two-Stage Stochastic MILP. Under this paradigm, hub activation ($y_i$) is decided in Stage~1 prior to the realization of hazard uncertainty, while second-stage dispatch ($x_{ijr}$) responds dynamically across discrete probabilistic weather scenarios characterized by varying flood extents and demand surges.

    \item \textbf{Integrated Multi-Commodity Allocation-Routing (MILP-VRP):}
    Future work should couple the proposed deprivation-minimizing macro-allocation model with a micro-level Vehicle Routing Problem with Time Windows (VRPTW). Incorporating dynamic road network damage graphs, travel time degradations, and heterogeneous vehicle fleet constraints (trucks, 4x4 vehicles, boats, and aerial drones) would bridge the gap between municipal quota determination and physical roadside delivery.

    \item \textbf{Dynamic Multi-Period Rolling Horizon Optimization:}
    Extending the model to a multi-period formulation ($t \in \{1, 2, \dots, T\}$) with state equations for inventory replenishment, rotting/spoilage rates, and non-linear cumulative deprivation accumulation would enable emergency managers to plan multi-day distribution schedules and track long-term recovery trajectories.

    \item \textbf{Empirical Epidemiological Calibration of Deprivation Functions:}
    Interdisciplinary research collaborating with public health researchers and emergency medical professionals should aim to empirically calibrate the deprivation parameters ($\lambda_r$) using local epidemiological health records, dehydration morbidity curves, and post-calamity rescue logs from recent Philippine typhoons and flash floods.

    \item \textbf{Geospatial Web-Based Decision Support System (DSS):}
    From a software engineering perspective, the underlying Pyomo optimization engine should be wrapped in an open-source, containerized web API integrated with a geographic information system (GIS) frontend. Developing an interactive, intuitive dashboard for the Davao City DRRMOC would enable operational duty officers to adjust scenario sliders, import real-time sensor data, solve the MILP in sub-second intervals, and visualize optimal dispatch routes on interactive maps.
\end{enumerate}
